\section{实验与结果}

\subsection{评测模型}

论文在三个制造任务上评估了 18 个 MLLM，覆盖开源与闭源模型。评测对象是通用基础多模态模型，而不是专门制造模型或专门 3D 模型。这一点与论文目标一致：它要回答“当前主流通用 MLLM 离真实制造认知还有多远”，而不是“专门模型能否把某个制造任务做高”。

\subsection{评测协议}

FORGE 中的任务被构造为多项选择题。对于 WORKVERI 和 ASSYVERI，每个样本包含 4--6 个部件。在图像评测中，每个选项对应图中某个部件，通常通过归一化中心坐标定位，例如“位于 $[0.70,0.44]$ 的部件”。在三视图评测中，作者采用 Set-of-Mark 视觉提示，用 A--F 等字母标注各个部件。

对于 SURFINSP，模型需要把工件表面状态分为五类之一：

\begin{itemize}
    \item crack
    \item cut
    \item deformation
    \item dent
    \item good
\end{itemize}

论文采用三种信息量逐步增加的评测设置：

\begin{itemize}
    \item Zero-Shot：只给测试图像或三视图和问题，不提供额外示例或参考。
    \item Reference-Conditioned：额外给出正确装配或无缺陷表面的参考图像，检验显式视觉参考是否帮助异常检测。
    \item In-Context Demonstration：在参考条件基础上加入完整求解示例，检验示范推理是否能弥补任务理解不足。
\end{itemize}

评价指标是 exact-match accuracy。模型自由文本回答中的选项字母会被抽取出来，与真实标签比较，预测完全一致即为正确。

\subsection{主要实验结论}

\subsubsection{语义理解强于形态分析}

当前 MLLM 在 WORKVERI 和 ASSYVERI 上整体好于 SURFINSP。这说明模型在宏观部件识别、类别判断和部分装配推理上已经有一定能力，但在表面缺陷、微观结构、局部形态变化等 morphology 任务上明显不足。

SURFINSP 看似只是判断缺陷类型，但它要求模型理解细微裂纹、凹痕、切损和变形等物理细节。这类信号对通用 MLLM 很不友好，因此成为三类任务中最困难的一项。

\subsubsection{领域知识不足是主要瓶颈}

在图像模态下，简单的 Reference-Conditioned 并不稳定提升 WORKVERI 和 ASSYVERI，甚至可能导致性能下降。相反，包含完整推理示范的 ICD 往往带来更普遍的提升。

这说明模型缺的不是“多看一张正常参考图”这么简单，而是对任务逻辑、制造规则和推理路径的理解。换言之，模型可能能看到零件，却不知道该如何基于规格判断它是否合格。

\subsubsection{三视图中的额外示例可能造成空间混淆}

在三视图模态下，Zero-Shot 反而常常优于 Ref-Cond 和 ICD。一个合理解释是：三视图依赖多角度空间表征，额外参考或示例没有被模型有效整合，反而引入空间对应混淆。

这说明当前通用 MLLM 对显式 3D 空间上下文仍然不够稳定。三视图虽然把点云转为模型可读的图像，但模型仍需要跨视角融合、部件对应和空间关系理解，这些能力并不等同于普通 2D 图像问答能力。

\subsubsection{型号级任务比工件级任务更难}

WORKVERI 和 ASSYVERI 中，模型在工件级任务上的表现普遍优于型号级任务。工件级任务只需区分大类，例如 screw、nut、washer；型号级任务则要求区分 M10/M12、长短差异、螺纹细节或同类零件中的规格差别。

这种差距说明，当前 MLLM 已具备一定制造对象理解能力，但对制造领域特异的细粒度差异仍然不足。

\subsection{定性错误案例}

论文的错误分析显示，MLLM 有时会尝试推断材料属性或服役状态，但这些推断并不总是可靠。一个案例中，模型错误判断某个零件材料，并把这个错误材料判断用于最终推理，表现出对材料属性的过度依赖。

另一个案例中，模型能识别工件类型，却误判型号或尺寸来源。例如它可能认为 Nut 太小，而真实问题是 Step Block 太大。尽管答案错误，模型中间推理中提到磨损、崩边等服役状态，也显示出未来用于预测性维护的潜力。

\subsection{瓶颈分析}

\subsubsection{Visual Grounding 不是主要瓶颈}

作者设计 probing 任务来单独测试视觉定位能力。单图定位要求模型在字母标签和坐标之间转换，例如 L$\rightarrow$C 或 C$\rightarrow$L；跨图对应要求模型在两张图中找到同一部件的对应位置。

结果显示，强模型在单图定位上接近天花板，例如部分模型在 L$\rightarrow$C 方向超过 97\%。这说明 WORKVERI 和 ASSYVERI 中的失败不能简单归因于“模型找不到部件”。主要问题更可能发生在定位之后的制造知识理解与规则推理。

\subsubsection{细粒度部件识别是领域知识瓶颈}

作者进一步设计缺失部件场景，给模型明确的装配规格、部件列表、数量和功能描述，再要求它判断缺少哪个部件。模型在螺丝、螺母、锚栓、楔块等视觉差异明显的部件上表现较好，但在 flat washer 等细粒度变体上系统性困难。

这说明模型常能发现“少了某个垫圈”，却不能准确判断“少的是哪一种垫圈”。由于 visual grounding 已被证明较强，这种错误更说明模型缺少不同零件变体之间功能差异和形态差异的制造知识。

\subsubsection{视觉投影是必要的}

论文测试了把点云坐标直接序列化成文本输入的方案。结果在 SURFINSP 上接近随机基线，说明文本坐标无法有效替代视觉渲染。只有在粗粒度 WORKVERI 上，坐标文本略高于随机，说明它可能提供一点形状分布信息，但远不足以支持细粒度制造理解。

因此，对没有原生 3D 编码器的通用 MLLM 来说，多视图视觉投影是处理 3D 制造数据的更有效接口。

\subsection{从 Benchmark 到训练资源}

FORGE 还被用作训练资源。作者用任务特定 SFT 微调 Qwen2.5-VL-3B-Instruct，并采用基于场景的训练/测试划分，以验证模型是否学到可迁移制造知识，而不是记住具体布局。

原始笔记记录：在 WORKVERI 三视图上，SFT 带来 90.8\% 的提升，使 3B 模型达到 Qwen3-VL-235B 的水平；在 ASSYVERI 图像上，SFT 带来 27.1\% 的相对增益，超过除 Gemini-3-Flash 和 GPT-5.2 之外的大多数模型。

这些提升发生在训练未见过的产品类别上，说明 FORGE 的细粒度标注确实编码了可迁移的制造知识。它不仅是考试卷，也是可操作的训练材料。
